\documentclass[11pt]{article}
\usepackage[utf8]{inputenc}
\usepackage[T1]{fontenc}
\usepackage[spanish,es-noshorthands]{babel}
\usepackage{lmodern}
\usepackage[margin=2.5cm]{geometry}
\usepackage{amsmath,amssymb,amsthm}
\usepackage{booktabs,tabularx,array}
\usepackage{xcolor}
\usepackage{enumitem}
\usepackage[colorlinks=true,linkcolor=blue!50!black,urlcolor=blue!50!black]{hyperref}

\setlist{itemsep=2pt,topsep=4pt}
\newcommand{\code}[1]{\texttt{\small #1}}
\newcommand{\act}{\mathrm{act}}
\newcommand{\term}{\mathrm{term}}
\newcommand{\none}{\mathrm{none}}
\newcommand{\used}{\mathrm{used}}
\newcommand{\new}{\mathrm{new}}
\newcommand{\buy}{b}
\newcommand{\1}{\mathbf{1}}
\newcommand{\Osc}{\mathcal{O}_{\mathrm{ch}}}
\newcommand{\Qbar}{\bar Q}

\newtheorem*{obs}{Observación}

\title{La verosimilitud de cada diseño del Monte Carlo\\[2pt]
\large con énfasis en la mezcla del diseño 3 ($r$ y chatarreo no observados)}
\author{Notas para Rodrigo Aldrette}
\date{15 de octubre de 2026}

\begin{document}
\maketitle

\noindent Versión en \LaTeX{} de \code{claude/niu/docs/loglikelihood.md}, revisada contra
\code{modelo\_tesis/ll\_estim.py} (\code{to\_cells}, \code{cell\_data}, \code{cell\_logp}).
Se añade la sección~\ref{sec:mezcla3}, que desarma qué entra nuevo en la mezcla cuando no se
observan ni la reparación ni el motivo de salida del coche.

\tableofcontents

%======================================================================
\section{Notación}
%======================================================================

\paragraph{Estados.} Dos espacios, como en Gillingham:
\begin{itemize}
  \item $\mathcal X$ (inicio de año): $\act(j,a,w)$, coche activo de marca $j$, edad
  $a=1,\dots,24$ y desgaste $w$; $\term(j)$, el coche se murió (accidente o fin de vida);
  y $\none$.
  \item $\mathcal H$ (tenencia después de comerciar): $\used(j,d,w)$, $\new(j)$ y $\none$.
\end{itemize}
Con $J=2$, $A=25$ y $W=71$ puntos de $w$: $n_{\act}=2\cdot 24\cdot 71=3{,}408$ y
$n=n_{\act}+J+1=3{,}411$ estados en $\mathcal X$ (y en $\mathcal H$).

\paragraph{Dentro de un año}, quien empieza en $x$:
\begin{enumerate}
  \item \textbf{Etapa 1} (escala $\sigma$): \emph{keep}, \emph{purge} (deshacerse y quedarse sin
  coche) o \emph{trade} (deshacerse y comprar). Al deshacerse de un activo elige vender
  ($\mu P - T_s$) o chatarrear (escala $\sigma_{\text{sell}}$). Con trade elige qué
  comprar: nuevo $j$ o usado $(j,d)$ y, dentro del nido, qué $w$ (escala $\sigma_w$).
  \item \textbf{Etapa 2} (escala $\sigma_{\text{rep}}$): con un usado $h$ de edad $1..23$
  decide reparar ($r=1$) o no ($r=0$). Nuevo, edad 24 y $\none$ no reparan.
  \item \textbf{Se maneja $h$}: se descompone con probabilidad $s(h)$ y pasa a $\term$; si
  sobrevive,
  \[
    w' = w + \delta_j - \text{acc\_age}_j - \kappa\, r + \eta,\qquad
    \eta\sim N(0,\sigma_\eta^2)\ \text{discretizado},
  \]
  y la edad pasa a $d+1$.
\end{enumerate}

\paragraph{Objetos del equilibrio} (los precios $P(j,a,w)$ nunca se observan; se resuelven
dentro de cada evaluación, régimen por régimen). Para el régimen $t$ y el tipo $\tau$:

\begin{center}
\begin{tabular}{lll}
\toprule
símbolo & qué es & código\\
\midrule
$C_{t\tau}(x,o)$ & $\Pr(\text{resultado } o\mid x)$ & \code{outcome\_probs}\\
$\buy_{t\tau}(h)$ & $\Pr(\text{comprar } h\mid \text{trade})$, incluye la elección de $w$ & \code{c.buy}\\
$p_{t\tau}(h)$ & $\Pr(r=1\mid h)$ & \code{c.repair}\\
$s(h)=s(j,d,w)$ & $\operatorname{logit}^{-1}(\text{acc\_int}_j+\text{acc\_age}_j\, d + w)$ & \code{breakdown\_prob}\\
$F_r(w'\mid j,w)$ & transición de $w$ dado $r$ & \code{w\_transition}\\
\bottomrule
\end{tabular}
\end{center}

\paragraph{Resultados observables.} $o\in\{0,\dots,4\}$: 0 keep, 1 purge-vende,
2 purge-chatarrea, 3 trade-vende, 4 trade-chatarrea. Escribo
$\Osc=\{2,4\}$ para ``chatarrear''. Un coche terminal siempre se chatarrea:
$C(\term(j),o)=0$ si $o\notin\Osc$.

\paragraph{Dos propiedades que hacen manejable la verosimilitud.}
\begin{enumerate}[label=(P\arabic*)]
  \item \label{P1} \textbf{$s$ no depende de $r$} (``opción 2'' de \code{modelo\_fin.md},
  sec.~3.3): reparar hoy no cambia la probabilidad de descomponerse hoy, solo el $w'$ de
  mañana.
  \item \label{P2} \textbf{$p(h)$ depende solo de $h$} (y de $t,\tau$), no de cómo se llegó
  a $h$: el que se quedó el coche y el que lo acaba de comprar reparan con la misma
  probabilidad.
\end{enumerate}

%======================================================================
\section{Qué es cada probabilidad y qué se observa}\label{sec:lectura}
%======================================================================

\subsection{La línea de tiempo de un hogar}

El año $k$ de un hogar tiene tres momentos, y en cada uno hay una variable:
\[
  \underbrace{h_{k-1}}_{\substack{\text{coche con el que}\\\text{terminó el año } k-1}}
  \xrightarrow[\text{se maneja}]{\;r_{k-1},\ \text{accidente},\ \eta\;}
  \underbrace{x_k}_{\substack{\text{estado de ese coche}\\\text{al empezar el año } k}}
  \xrightarrow[\text{etapa 1}]{\;o_k\;}
  \underbrace{h_k}_{\substack{\text{coche con el que}\\\text{termina el año } k}}
  \xrightarrow[\text{etapa 2}]{\;r_k\;}\cdots
\]
\begin{itemize}
  \item $h_{k-1}$: qué coche tenía el hogar después de comerciar el año pasado, p.~ej.
  $\used(\text{LB},5,w=0.30)$.
  \item $x_k$: qué le pasó a \emph{ese mismo coche} al manejarlo un año. O se murió
  ($\term(\text{LB})$) o sigue vivo un año más viejo y con desgaste nuevo
  ($\act(\text{LB},6,w')$). El hogar todavía no ha decidido nada.
  \item $o_k$: qué hace el hogar con ese coche (etapa 1).
  \item $h_k$: con qué coche se queda al final. Si $o_k=$ keep, $h_k$ es el mismo coche
  ($\used(\text{LB},6,w')$); si purge, $h_k=\none$; si trade, $h_k$ es el que compró.
\end{itemize}

\subsection{Las tres piezas de probabilidad}

\paragraph{$Q_r(h_{k-1},x_k)$: física, no es decisión.}
$\Pr(x_k=x\mid h_{k-1}=h,\ r_{k-1}=r)$: la probabilidad de que el coche $h$ amanezca en el
estado $x$, dado que el año pasado se reparó ($r=1$) o no. Ver~\eqref{eq:Qr}.

\paragraph{$C_t(x,o)$: decisión de etapa 1, sin el qué comprar.}
\[
  C_t(x,o)=\Pr(o_k=o\mid x_k=x,\ \text{régimen } t,\ \text{tipo }\tau).
\]
Es una matriz de $n\times 5$: una fila por estado de inicio de año, una columna por
resultado. Se arma con dos logits anidados. Sea $\text{keep}(x)$, $\text{purge}(x)$,
$\text{trade}(x)$ la elección de etapa 1 (escala $\sigma$) y
$\text{ch}(x)=\Pr(\text{chatarrear}\mid\text{se deshace del coche},x)$ la elección vender
vs chatarrear (escala $\sigma_{\text{sell}}$). Entonces (\code{outcome\_probs}):
\begin{center}
\begin{tabular}{cll}
\toprule
$o$ & nombre & $C_t(x,o)$\\
\midrule
0 & keep & $\text{keep}(x)$\\
1 & purge-vende & $\text{purge}(x)\,(1-\text{ch}(x))$\\
2 & purge-chatarrea & $\text{purge}(x)\,\text{ch}(x)$\\
3 & trade-vende & $\text{trade}(x)\,(1-\text{ch}(x))$\\
4 & trade-chatarrea & $\text{trade}(x)\,\text{ch}(x)$\\
\bottomrule
\end{tabular}
\end{center}
Casos: si $x=\term(j)$, no hay keep y $\text{ch}=1$ (el coche muerto solo se chatarrea),
así que solo $o\in\{2,4\}$ tienen probabilidad. Si $x=\none$, $o=0$ es ``seguir sin coche''
y $o=3$ es ``comprar'' (no hay nada que vender ni chatarrear). Cada fila suma 1.

\paragraph{$\buy_t(h_k)$: qué se compra, si hubo trade.}
$\Pr(h_k=h\mid o_k\in\{3,4\})$: marca, edad y $w$ del coche comprado (nested logit con
$\sigma_{\text{trade}}$ y $\sigma_w$). No depende de $x$: lo que se recibe por el coche viejo
entra igual en todas las opciones de compra.

\paragraph{$p_t(h_k)$: decisión de etapa 2.} $\Pr(r_k=1\mid h_k)$. Afecta el
$Q_{r}(h_k,x_{k+1})$ \emph{del año siguiente}.

\paragraph{Juntando.} Dado $x_k$, el par $(o_k,h_k)$ tiene probabilidad
$C_t(x_k,o_k)\cdot\buy_t(h_k)^{\1\{o_k\in\{3,4\}\}}$; si no hubo trade, $h_k$ queda
determinado por $(x_k,o_k)$ y no aporta nada más. La probabilidad de una fila completa es
\[
  \Pr(x_k,o_k,h_k\mid h_{k-1})=
  \underbrace{\sum_{r}\pi_t(r\mid h_{k-1})\,Q_r(h_{k-1},x_k)}_{\text{física, con o sin mezcla}}
  \cdot\underbrace{C_t(x_k,o_k)}_{\text{etapa 1}}
  \cdot\underbrace{\buy_t(h_k)^{\text{trade}}}_{\text{qué compra}}.
\]

\paragraph{Ejemplo.} $h_{k-1}=\used(\text{LB},5,0.30)$, $x_k=\act(\text{LB},6,0.10)$,
$o_k=$ trade-vende, $h_k=\used(\text{HB},3,-0.15)$. La probabilidad (diseño 2) es
\[
  \underbrace{(1-s)\bigl[(1-p)F_0(0.10\mid 0.30)+pF_1(0.10\mid 0.30)\bigr]}_{\text{sobrevivió y llegó a } w'=0.10}
  \times\underbrace{\text{trade}(x_k)(1-\text{ch}(x_k))}_{C(x_k,3)}
  \times\underbrace{\buy(\used(\text{HB},3,-0.15))}_{\text{eligió ese coche}}.
\]

\subsection{Qué se observa en cada diseño, comparado con Gillingham}

Las tres variables de cada año son $x_k$, $o_k$, $h_k$. $h_k$ y $o_k$ se observan siempre
(registro de propiedad y si el coche viejo se vendió o se chatarreó). Lo que cambia es $x_k$
y $r$:

\begin{center}
\small
\begin{tabularx}{\textwidth}{l>{\raggedright\arraybackslash}X>{\raggedright\arraybackslash}X>{\raggedright\arraybackslash}X}
\toprule
 & $x_k$ si el coche sigue (keep o vende) & $x_k$ si el coche se chatarrea & $r$\\
\midrule
Gillingham (registro danés) & implícito: $\act(j,a+1)$, no hay $w$ & no: no se sabe si murió o lo chatarrearon & no existe\\
Diseño 1 & sí, con $w'$ & sí: si murió y, si no, su $w'$ & sí\\
Diseño 2 & sí, con $w'$ & sí & no\\
Diseño 3 & sí, con $w'$ & no & no\\
\bottomrule
\end{tabularx}
\end{center}

\noindent Es decir, el diseño 3 es Gillingham más una medición de $w$ (p.~ej. inspección u
odómetro) en los coches que siguen circulando. En Gillingham, sin $w$, para un coche vivo
$x_k$ es una función determinista de $h_{k-1}$, así que la única incógnita en $x_k$ es
``¿se murió?''; por eso basta con $h_{k-1}\to(o_k,h_k)$. Aquí, con $w'$ aleatorio, $x_k$
lleva información propia (el $w'$) y es justo la que usa Hu \& Xin.

\subsection{¿Hay que seguir hogares o bastan las celdas?}

Bastan los pares consecutivos $(h_{k-1};\,x_k,o_k,h_k)$; el panel solo sirve para
\emph{formar} esos pares. La razón es que la verosimilitud de un hogar se factoriza:
\[
  \Pr(\text{historia del hogar}\mid h_0)=\prod_{k=1}^{K-1}\Pr(x_k,o_k,h_k\mid h_{k-1}),
\]
y cada $r_{k-1}$ no observado se integra \emph{dentro de su propio factor}. Esto vale porque:
\begin{enumerate}
  \item $r_{k-1}$ solo afecta $x_k$ (el desgaste de un año); no afecta $x_{k+1}$ más que a
  través de $h_k$, que se observa.
  \item $r_{k-1}$ se elige con choques de valor extremo independientes entre años; dado
  $h_{k-1}$, no está correlacionado con $r_{k-2}$ ni con nada del pasado.
  \item No hay heterogeneidad no observada persistente: el tipo de hogar y $w$ se observan.
\end{enumerate}
Por eso no hace falta sumar sobre secuencias $(r_0,r_1,r_2)$ de cada hogar
($2^{K-1}$ combinaciones): la suma se hace año por año. Y como la probabilidad de cada par
depende solo de $(t,\tau,h_{k-1},x_k,o_k,h_k)$, juntar pares idénticos en celdas con
conteo es exacto.

La identificación de $r$ tampoco viene de seguir al hogar, sino de comparar la distribución
de $w'$ dado el mismo $h$ entre regímenes (ec.~\eqref{eq:d2}): con $R_t$ más barato hay más
peso en $F_1$. Seguir hogares ayudaría si hubiera un tipo no observado (un hogar que repara
siempre); entonces la mezcla sería sobre el tipo, la historia completa del hogar entraría
en un mismo término y las celdas por par ya no bastarían.

En el diseño 3 la factorización también se mantiene: aunque $x_k$ no se vea en una salida,
el factor siguiente condiciona en $h_k$, que sí se ve.

%======================================================================
\section{Los datos: panel y celdas}
%======================================================================

\subsection{El panel (igual en los tres diseños)}

\code{gen\_dataset.simulate\_panel}: por régimen $t\in\{\zeta=-0.3,0,+0.3\}$, $N=10{,}000$
hogares repartidos entre tipos, $K=4$ años, con estado inicial sacado de la distribución
estacionaria del régimen. Una fila por hogar y año con
$(\text{régimen},\text{tipo},\text{id},k,x,o,h,r)$, donde $r=-1$ si $h$ no podía repararse.

Siempre se observa: régimen, tipo, $x$ (incluido $w$), $o$, $h$ (incluido el $w$ del usado
comprado). Nunca: precios de usados ni choques de utilidad.

\begin{center}
\begin{tabular}{cclllr}
\toprule
id & $k$ & $x$ & $o$ & $h$ & $r$\\
\midrule
17 & 0 & $\act(\text{LB},5,0.30)$ & keep & $\used(\text{LB},5,0.30)$ & 1\\
17 & 1 & $\act(\text{LB},6,0.10)$ & keep & $\used(\text{LB},6,0.10)$ & 0\\
17 & 2 & $\term(\text{LB})$ & trade-chatarrea & $\used(\text{HB},3,-0.15)$ & 0\\
17 & 3 & $\act(\text{HB},4,0.05)$ & purge-vende & $\none$ & $-1$\\
\bottomrule
\end{tabular}

\smallskip
{\small Ejemplo ilustrativo de un hogar.}
\end{center}

\subsection{La unidad de observación: $h_{k-1}\to(x_k,o_k,h_k)$}

La verosimilitud es \textbf{condicional en la tenencia del año anterior}. Cada fila con
$k\ge 1$ aporta
\[
  \Pr(x_k,o_k,h_k\mid h_{k-1},t,\tau).
\]
El año $k=0$ solo da $h_0$; no se usa la distribución inicial. Por réplica:
$3\times 10{,}000\times 3 = 90{,}000$ observaciones. En el diseño~1 la transición además
lleva $r_{k-1}$.

\subsection{Celdas con conteos}

Todo es discreto, así que las observaciones idénticas se agrupan (\code{to\_cells}):
\[
  \ell(\theta)=\sum_{c}\text{cnt}_c\,\log\Pr\bigl(c\mid\theta,P(\theta)\bigr).
\]
Llaves: $(t,\tau,h_{\text{prev}},r_{\text{prev}},x,o,h)$. Cada diseño borra información
\emph{antes} de agrupar:

\begin{center}
\begin{tabular}{clll}
\toprule
diseño & $r_{\text{prev}}$ & $x$ & celdas/réplica\\
\midrule
1 & observado & observado & 48,729\\
2 & borrado ($-1$) & observado & 45,802\\
3 & borrado & borrado ($-1$) si $h_{\text{prev}}\neq\none$ y $o\in\Osc$ & 45,273\\
\bottomrule
\end{tabular}
\end{center}

Agrupar es exacto: la probabilidad de una observación depende solo de sus llaves.

%======================================================================
\section{Bloques comunes}
%======================================================================

Fijo $h=h_{k-1}=\used(j,d,w)$ (o $\new(j)$, con $d=0$, $w=w_0$, $p=0$) y omito $t,\tau$
cuando no hay confusión.

\paragraph{Transición física condicional en $r$.} Defino la matriz
$Q_r:\mathcal H\to\mathcal X$:
\begin{equation}\label{eq:Qr}
  Q_r(h,x)=
  \begin{cases}
    s(h) & x=\term(j),\\
    \bigl(1-s(h)\bigr)\,F_r(w'\mid j,w) & x=\act(j,d+1,w'),\\
    0 & \text{otro caso,}
  \end{cases}
\end{equation}
con $s=1$ si $d=24$ y $Q_r(\none,\none)=1$.

\paragraph{Decisión del año $k$.}
\[
  \text{choice}(x,o,h_k)=C(x,o)\cdot \buy(h_k)^{\1\{o\ \text{es trade}\}}.
\]

\paragraph{Mezcla sobre $r$.} El objeto que reaparece en los diseños 2 y 3 es la transición
física \emph{promediada} con los pesos de reparación del régimen:
\begin{equation}\label{eq:Qbar}
  \boxed{\;\Qbar_t(h,x)=\bigl(1-p_t(h)\bigr)\,Q_0(h,x)+p_t(h)\,Q_1(h,x).\;}
\end{equation}
Por~\ref{P1}, $\Qbar_t(h,\term(j))=s(h)$ no depende de $p_t$: la mezcla solo actúa sobre
la parte de supervivencia.

%======================================================================
\section{Diseño 1: se observa $r$ y el motivo de salida}
%======================================================================

\[
  \Pr(r,x,o,h_k\mid h)=\underbrace{p(h)^r\bigl(1-p(h)\bigr)^{1-r}}_{\text{etapa 2}}
  \cdot \underbrace{Q_r(h,x)}_{\text{física}}\cdot \underbrace{C(x,o)\,\buy(h_k)^{\text{trade}}}_{\text{etapa 1}},
\]
con el primer factor igual a 1 si $h$ no puede repararse. Logit dinámico estándar con dos
decisiones observadas y transición observada. Sin mezcla. El log se separa en tres sumandos:
\begin{itemize}
  \item $\log\Pr(r\mid h)$: $\sigma_{\text{rep}}$ y la utilidad de reparar (vía la variación
  de $R_t$).
  \item $\log Q_r(h,x)$: $\kappa$, $\delta$, $\sigma_\eta$ (transiciones de $w$ dado $r$) y
  acc\_int, acc\_age (accidentes).
  \item $\log C+\log \buy$: $\mu,u_0,u_1,T_b,T_s,\sigma_{\text{sell}}$, como en Gillingham
  pero con $w$.
\end{itemize}
Hu \& Xin no se usa: es la referencia ``choices observed'' (su sec.~5).

%======================================================================
\section{Diseño 2: no se observa $r$}
%======================================================================

Se ve $x$ (con $w'$), $o$, $h$; no $r$. Se integra $r$:
\begin{equation}\label{eq:d2}
  \Pr(x,o,h_k\mid h)=\Qbar_t(h,x)\cdot C_t(x,o)\cdot\buy_t(h_k)^{\text{trade}}.
\end{equation}
Para un coche que sobrevive,
\[
  \Qbar_t\bigl(h,\act(j,d+1,w')\bigr)=\bigl(1-s(h)\bigr)\Bigl[(1-p_t(h))F_0(w'\mid j,w)+p_t(h)F_1(w'\mid j,w)\Bigr],
\]
y para un accidente, $s(h)$. Código: \code{trans = (1 - p1)*px[0] + p1*px[1]}.

\paragraph{Lo que se toma de Hu \& Xin.}
\begin{enumerate}
  \item \emph{Identificación} (secs.~3 y 6.3).
  \begin{itemize}
    \item Assumption 2 (transición aditiva): $w'=g(w,r)+\eta$, $\eta\perp r$, con
    $g(w,r)=w+\delta_j-\text{acc\_age}_j-\kappa r$.
    \item Assumption 7 (variable excluida): $R_t$ mueve $p_t(h)$ pero no $F_r$. La variación
    de la mezcla entre regímenes con componentes fijas separa $F_0,F_1,p_t$ (ecs.~3.2--3.4).
    Bastan 2 valores; hay 3.
    \item Observabilidad parcial (ec.~6.4): $f(x'\mid x,a^1)=\sum_{a^2}f(x'\mid x,a)\,
    p(a^2\mid x,a^1)$; aquí $p(a^2\mid x,a^1)=p(h)$ por~\ref{P2}.
  \end{itemize}
  \item \emph{Estimador} (sec.~4, ec.~4.1): MV con la decisión no observada sumada dentro
  del logaritmo, en su caso paramétrico.
\end{enumerate}

\paragraph{Lo que no se toma.} Sin sieve ($g$ lineal en $r$, $\eta$ normal); sin primer
paso no paramétrico (Remark~2; no implementado); y equilibrio: los CCPs dependen de precios
no observados que se resuelven en cada evaluación (punto fijo anidado), cosa que el paper,
de agente único, no cubre. La identificación de la mezcla no cambia porque los precios son
función de $(\theta,t)$, no de $r$.

\paragraph{Clave para lo que sigue.} En~\eqref{eq:d2} $C_t(x,o)$ \emph{multiplica fuera} de
la mezcla: como $x$ (y por tanto $w'$) se observa, $C_t(x,o)$ no depende de $r$ y se
factoriza. Además, como $x$ se observa en \emph{todas} las filas, el marginal
$\Pr(x\mid h)=\Qbar_t(h,x)$ es un objeto poblacional directamente observable: es literalmente
la mezcla finita de Hu \& Xin.

%======================================================================
\section{Diseño 3: no se observa $r$ ni el motivo de salida}
%======================================================================

Como el diseño 2, salvo cuando un coche sale del parque: si $h_{k-1}\ne\none$ y
$o\in\Osc$, no se sabe si fue
\begin{itemize}
  \item un \textbf{accidente}: $x=\term(j)$, chatarreo forzoso; o
  \item un \textbf{chatarreo voluntario}: $x=\act(j,d+1,w')$ con $w'$ no observado y el dueño
  eligió chatarrear en vez de vender.
\end{itemize}
En esas filas $x$ se borra (\code{d.loc[salida, "x"] = -1}). Es lo que ve Gillingham en el
registro danés (apéndice D).

Filas sin salida: igual que~\eqref{eq:d2}. Filas con salida se suma sobre $r$ \emph{y} sobre
$x$:
\begin{multline}\label{eq:d3}
  \Pr(o,h_k\mid h)=\buy_t(h_k)^{\text{trade}}\sum_{r\in\{0,1\}}\pi_t(r\mid h)
  \Bigl[s(h)\,C_t(\term(j),o)\Bigr.\\
  \Bigl.+\bigl(1-s(h)\bigr)\sum_{w'}F_r(w'\mid j,w)\,
  C_t\bigl(\act(j,d+1,w'),o\bigr)\Bigr],
\end{multline}
con $\pi_t(1\mid h)=p_t(h)$, $\pi_t(0\mid h)=1-p_t(h)$. Código: \code{mix}, \code{lp\_exit} en
la rama \code{design == 3} de \code{cell\_logp}.

%======================================================================
\section{Qué añade la mezcla en el diseño 3}\label{sec:mezcla3}
%======================================================================

\subsection{Forma compacta: la verosimilitud parcial de Gillingham con $\Qbar_t$}

Como desde $h$ solo se puede llegar a $\term(j)$ o a $\act(j,d+1,\cdot)$, la suma sobre $x$
en~\eqref{eq:d3} es un producto de matrices:
\begin{equation}\label{eq:compacta}
  \boxed{\;\Pr(o,h_k\mid h)=\bigl(\Qbar_t\,C_t\bigr)(h,o)\cdot\buy_t(h_k)^{\text{trade}},\qquad o\in\Osc.\;}
\end{equation}
Comparar con Gillingham (sec.~\ref{sec:gill}): su verosimilitud parcial es
$(QC)(h,o)\,\buy^{\text{trade}}$. El diseño 3 es exactamente esa expresión con la matriz
física sustituida por la mezcla $\Qbar_t$ de~\eqref{eq:Qbar}. Del mismo modo, el diseño 2 es
la ``completa'' de Gillingham, $Q(h,x)C(x,o)$, con $Q\to\Qbar_t$. Es decir:

\begin{center}
\begin{tabular}{lcc}
\toprule
 & $x$ observado & $x$ no observado en salidas\\
\midrule
$r$ observado & $Q_r(h,x)\,C(x,o)$ \ (diseño 1) & $(Q_rC)(h,o)$\\
$r$ no observado & $\Qbar_t(h,x)\,C(x,o)$ \ (diseño 2) & $(\Qbar_tC)(h,o)$ \ (diseño 3)\\
\bottomrule
\end{tabular}
\end{center}

\noindent (La celda $(Q_rC)$ no es un diseño del MC; está para completar la tabla.)

\subsection{Desarmando la mezcla: dos variables latentes anidadas}

Defino la probabilidad esperada de chatarrear con resultado $o$, dado que el coche
sobrevive y se eligió $r$:
\begin{equation}\label{eq:Cbar}
  \bar C^{(r)}_t(h,o)=\sum_{w'}F_r(w'\mid j,w)\,C_t\bigl(\act(j,d+1,w'),o\bigr).
\end{equation}
Usando $\sum_r\pi_t(r\mid h)=1$, la ecuación~\eqref{eq:d3} queda
\begin{equation}\label{eq:d3b}
  \frac{\Pr(o,h_k\mid h)}{\buy_t(h_k)^{\text{trade}}}=
  \underbrace{s(h)\,C_t(\term(j),o)}_{\text{accidente}}
  +\bigl(1-s(h)\bigr)\underbrace{\Bigl[(1-p_t(h))\,\bar C^{(0)}_t(h,o)+p_t(h)\,\bar C^{(1)}_t(h,o)\Bigr]}_{\text{chatarreo voluntario, mezclado sobre } r}.
\end{equation}

Hay tres cosas nuevas respecto al diseño 2.

\paragraph{(i) Una segunda variable latente: el estado $x$.} En el diseño 2 la única latente
es $r$. Aquí también lo es el estado de inicio de año del coche que sale: muerto
($\term$) o vivo con algún $w'$. Esta mezcla es la del apéndice D de Gillingham, no la de
Hu \& Xin. Su peso, $s(h)$, depende de $w$ (observado en $h$) y no de $r$.

\paragraph{(ii) Las dos mezclas no se cruzan gracias a~\ref{P1}.} El término de accidente
sale de la suma sobre $r$ porque $\Qbar_t(h,\term)=s(h)$ no depende de $p_t$. La mezcla
sobre $r$ queda \emph{anidada dentro} de la rama ``sobrevivió''. Con la opción~1
(\code{opcion\_1\_reparar.md}), donde reparar baja la probabilidad de descomponerse ya este
año, el accidente sería $\sum_r\pi_t(r\mid h)\,s_r(h)\,C_t(\term,o)$: la reparación movería
los pesos accidente/voluntario y las dos mezclas quedarían entrelazadas (y con selección por
supervivencia en los pesos).

\paragraph{(iii) $C$ entra \emph{dentro} de la mezcla sobre $r$.} Es la diferencia de fondo.
En el diseño 2, $C_t(x,o)$ se factoriza porque $w'$ se observa. En el diseño 3, la
probabilidad de chatarrear voluntariamente depende de $w'$ (un coche más gastado vale menos
en el mercado secundario, así que venderlo rinde menos frente a chatarrearlo), y $w'$ depende
de $r$. Entonces las ``componentes'' de la mezcla ya no son las densidades $F_r$ sino sus
\emph{momentos} contra $C_t$:
\[
  F_r(\cdot\mid j,w)\ \longrightarrow\ \bar C^{(r)}_t(h,o)=\mathbb E_{F_r}\bigl[C_t(\act(j,d+1,w'),o)\bigr].
\]
Se pierde la densidad de $w'$ y queda un escalar por $(h,o,r)$.

\subsection{Cuánto informa una salida sobre la reparación}

La sensibilidad de la probabilidad de una salida al peso de reparación es
\begin{equation}\label{eq:dp}
\begin{aligned}
  \frac{\partial}{\partial p_t(h)}\,\frac{\Pr(o,h_k\mid h)}{\buy^{\text{trade}}}
  &=\bigl(1-s(h)\bigr)\Bigl[\bar C^{(1)}_t(h,o)-\bar C^{(0)}_t(h,o)\Bigr]\\
  &=\bigl(1-s(h)\bigr)\sum_{w'}\bigl[F_1-F_0\bigr](w'\mid j,w)\;C_t\bigl(\act(j,d+1,w'),o\bigr).
\end{aligned}
\end{equation}
Lecturas:
\begin{itemize}
  \item Si $\kappa>0$, $F_1$ es $F_0$ desplazada hacia $w'$ menores (menos desgaste). Si
  $C_t(\act(\cdot,w'),o)$ crece con $w'$ para $o\in\Osc$, entonces~\eqref{eq:dp} es
  negativa: donde más se repara, menos chatarreo voluntario.
  \item La información es de \emph{primer orden en $\kappa$ por la pendiente de $C$ en $w'$}:
  si el chatarreo voluntario es casi plano en $w'$ (p.~ej. $\sigma_{\text{sell}}$ grande o
  precios poco sensibles al desgaste), las salidas casi no informan sobre $(p,\kappa)$.
  \item Se multiplica por $(1-s(h))$: en coches viejos o muy gastados, donde casi toda
  salida es accidente, la salida no dice nada sobre la reparación.
\end{itemize}

\subsection{El precio: $w'$ faltante no al azar}

En el diseño 3 el $w'$ de un coche vivo se observa solo si el dueño \emph{no} lo chatarrea.
Lo que el econometrista ve para un $h$ dado es
\[
  \Pr\bigl(x=\act(j,d+1,w'),\,o\mid h\bigr)=\bigl(1-s(h)\bigr)\sum_r\pi_t(r\mid h)F_r(w'\mid j,w)\cdot C_t(\act(j,d+1,w'),o),
  \qquad o\notin\Osc,
\]
más la masa total de salidas~\eqref{eq:d3b}. Sumando sobre $o\notin\Osc$, la densidad de
$w'$ observada es
\[
  \bigl(1-s(h)\bigr)\,\Bigl[\sum_r\pi_t(r\mid h)F_r(w'\mid j,w)\Bigr]\cdot\Bigl(1-C_t^{\text{ch}}(\act(j,d+1,w'))\Bigr),
  \qquad C^{\text{ch}}_t=\sum_{o\in\Osc}C_t(\cdot,o),
\]
es decir, la mezcla de Hu \& Xin \emph{inclinada} por la probabilidad de no chatarrear, que
depende del propio $w'$ (faltante no al azar). Dos consecuencias:

\begin{obs}
\begin{enumerate}
  \item \textbf{La verosimilitud sigue siendo correcta.} $C_t$ es un objeto del modelo, así que
  la inclinación se corrige dentro de~\eqref{eq:d2} y~\eqref{eq:d3}: no hay sesgo de
  selección en el MLE si el modelo está bien especificado.
  \item \textbf{El argumento de identificación de Hu \& Xin ya no se aplica tal cual.} Su
  Assumption~7 pide componentes de la mezcla invariantes entre regímenes. En el diseño 3 las
  componentes observables son $F_r(w')\,(1-C^{\text{ch}}_t(w'))$, que sí cambian con $t$
  (el precio de reparación cambia el valor de los coches y con ello la decisión de
  chatarrear). El marginal limpio $\Qbar_t(h,\cdot)$ del diseño 2 deja de ser un objeto
  observable. La separación de $F_0$, $F_1$ y $p_t$ se apoya ahora también en la estructura
  de $C_t$ (el modelo de etapa 1), no solo en la exclusión de $R_t$.
\end{enumerate}
\end{obs}

En la práctica la pérdida es moderada si la tasa de chatarreo voluntario es baja (la
inclinación $1-C^{\text{ch}}_t$ es casi constante en $w'$), y crece en edades altas, donde
el chatarreo voluntario es frecuente. Es la contraparte de~\eqref{eq:dp}: cuanto más
depende el chatarreo de $w'$, más informa la salida sobre la reparación, pero también más
se distorsiona la muestra de $w'$ observados.

\subsection{Casos especiales en el código}

\begin{itemize}
  \item $h=\new(j)$: $p=0$, no hay mezcla sobre $r$, pero sí la de accidente/voluntario.
  \item $h=\used(j,24,w)$: $s=1$, la salida es siempre accidente y~\eqref{eq:d3b} se reduce
  a $C_t(\term(j),o)$.
  \item $h=\none$: no hay coche que salga; $x=\none$ se observa siempre.
  \item En el diseño 3 $x=\term$ nunca aparece observado: toda fila con $x=\term$ tiene
  $o\in\Osc$ y se borra.
\end{itemize}

%======================================================================
\section{Para comparar: Gillingham sobre los mismos datos}\label{sec:gill}
%======================================================================

\code{gillingham/ll\_estim.py}, sobre el mismo panel (\code{montecarlo.to\_gillingham}): se
borran $w$ y $r$ y se juntan los regímenes. Celdas $(\tau,h_{\text{prev}},o,h)$ (parcial,
$\approx 1{,}065$) o $(\tau,h_{\text{prev}},x,o,h)$ (completa, $\approx 1{,}108$).
\begin{itemize}
  \item Parcial (la del paper, análoga al diseño 3): $(QC)(h,o)\,\buy^{\text{trade}}$.
  \item Completa (análoga al diseño 2): $Q(h,x)\,C(x,o)\,\buy^{\text{trade}}$.
\end{itemize}
Sin $w$ ni reparación el modelo está mal especificado y sus estimaciones están sesgadas
(\code{reporte\_mc.md}, sec.~2).

%======================================================================
\section{Cómo se maximiza (igual en los tres diseños)}
%======================================================================

\begin{itemize}
  \item \textbf{Punto fijo anidado.} Cada evaluación de $\ell(\theta)$ resuelve
  $z_t(\theta)=(\text{EV por tipo},P)$ de los 3 regímenes con pasos de cuerda desde el
  equilibrio anterior y predictor $z+(dz/d\theta)\Delta\theta$ (\code{LLEval},
  \code{equilibrio.solve\_chord}).
  \item \textbf{Gradiente implícito.} $dz_t/d\theta=-F_z^{-1}F_\theta$ (\code{dz\_dx}); el
  score incluye el efecto directo de $\theta$ y el indirecto vía precios.
  \item \textbf{Optimizador.} L-BFGS y luego BHHH desde la mejor evaluación, hasta
  $g'B^{-1}g/N<10^{-9}$.
  \item \textbf{Errores estándar.} Inversa de $\sum_c\text{cnt}_c\,s_cs_c'$, con método delta
  a la escala natural ($\mu,\kappa,\sigma_\eta,\sigma_{\text{rep}}$ en log;
  $\sigma_{\text{sell}}$ con sigmoide).
  \item \textbf{Parámetros libres:} 27. Fijos: $\sigma,\sigma_{\text{trade}},\sigma_w,\beta,
  p_{\text{new}},p_{\text{scrap}},u_2$ y $R$ (dato).
\end{itemize}

%======================================================================
\section{Resumen}
%======================================================================

\begin{center}
\small
\begin{tabularx}{\textwidth}{>{\raggedright\arraybackslash}p{3.3cm}
  >{\raggedright\arraybackslash}X >{\raggedright\arraybackslash}X >{\raggedright\arraybackslash}X}
\toprule
 & diseño 1 & diseño 2 & diseño 3\\
\midrule
observa $r$ & sí & no & no\\
observa $x$ en salidas & sí & sí & no\\
probabilidad & $\pi(r\mid h)\,Q_r(h,x)\,C(x,o)$ & $\Qbar_t(h,x)\,C(x,o)$ & salidas: $(\Qbar_tC)(h,o)$\\
mezcla sobre $r$ & no & sí, en $w'$; $C$ fuera & sí, en $w'$ y en el chatarreo; $C$ dentro\\
mezcla accidente/voluntario & no & no & sí (anidada por~\ref{P1})\\
$w'$ & siempre observado & siempre observado & faltante no al azar en salidas\\
papel de Hu \& Xin & referencia & identificación (A2, A7, ec.~6.4) y estimador (ec.~4.1) & mezcla sobre $r$ como en diseño 2; la identificación se apoya además en $C_t$; + apéndice D de Gillingham\\
\bottomrule
\end{tabularx}
\end{center}

\end{document}
